% =========================================================================
% Amazon Research Awards — AWS Cryptography, Fall 2026
% Submission window: 1 Oct – 4 Nov 2026 (23:59 PT). Decisions: Feb 2027.
% Max 4 pages excluding appendices — body is 4pp, measurements in Appendix A/B.
%
% Every quantitative figure is a measured elapsed time or an instrumented
% counter from a real run, taken from docs/experiments.md. Nothing is composed
% from a cost model. Claims about HEIR are sourced to arXiv:2508.11095.
% =========================================================================
\documentclass[10pt]{article}

\usepackage[utf8]{inputenc}
\usepackage[T1]{fontenc}
\usepackage{lmodern}
\usepackage[margin=0.88in]{geometry}
\usepackage{amsmath,amssymb}
\usepackage{booktabs}
\usepackage{array}
\usepackage{microtype}
\usepackage{xcolor}
\usepackage{enumitem}
\usepackage{titlesec}
\usepackage[hidelinks]{hyperref}

\newcommand{\TODO}[1]{\textcolor{red}{\textbf{[FILL: #1]}}}

\titlespacing*{\section}{0pt}{1.1ex plus .2ex}{0.5ex}
\titlespacing*{\subsection}{0pt}{1.1ex plus .2ex}{0.45ex}
\titleformat{\section}{\normalfont\large\bfseries}{\thesection}{0.6em}{}
\titleformat{\subsection}{\normalfont\normalsize\bfseries}{\thesubsection}{0.6em}{}
\setlist{nosep,leftmargin=1.3em}
\setlength{\parskip}{0.15ex}
\linespread{0.97}\selectfont
\pagestyle{plain}

\newcommand{\baru}{\textsc{Baru}}
\newcommand{\heir}{\textsc{Heir}}

\begin{document}

% -------------------------------------------------------------------------
\begin{center}
{\Large\bfseries Compiling Across the FHE/MPC Boundary}\\[0.3ex]
{\large A continuation-based compiler architecture for programs whose cheapest backend changes mid-computation}\\[1.0ex]
{\normalsize
\textbf{Principal Investigator:} \TODO{PI full name}, \TODO{title}\\
Departamento de Engenharia de Computação e Sistemas Digitais (PCS),
Escola Politécnica, USP --- LARC\\
\TODO{PI institutional e-mail} \quad\textbullet\quad \TODO{PI ORCID / Lattes}\\[0.5ex]
\textbf{Co-investigator:} Eduardo Morais, \TODO{role and affiliation}\\[0.5ex]
\textbf{Topic:} \emph{Compilers and development tools for FHE and MPC}
\quad\textbullet\quad
\textbf{Secondary:} \emph{Privacy-preserving AI inference and training on encrypted or distributed data}\\[0.4ex]
\textbf{Artifact:} \url{https://github.com/emmorais/baru} --- open source, Apache-2.0
}
\end{center}
\vspace{0.3ex}\hrule\vspace{1.0ex}

% -------------------------------------------------------------------------
\section{The architectural gap}

\heir{}~\cite{heir} is, by its own account and by community adoption, the de
facto compiler infrastructure for homomorphic encryption: an MLIR-based stack
with dialects spanning \texttt{secret}, \texttt{bgv}, \texttt{ckks},
\texttt{cggi}, library APIs for OpenFHE, Lattigo, TFHE-rs and Jaxite, and
hardware targets. We take that as given. \textbf{This proposal is not an attempt
to produce faster homomorphic code than \heir{}; inside \heir{}'s domain that
would be an engineering result, not a scientific one.}

What interests us is the boundary of that domain, which \heir{}'s own paper
draws precisely. Three consequences of a static, FHE-only lowering pipeline:

\begin{itemize}
\item \textbf{Secret-dependent control flow is eliminated, not executed.}
\texttt{convert-if-to-select} rewrites a secret-dependent branch into
\texttt{arith.select} --- so both arms are evaluated. Loops with
secret-dependent bounds are unsupported; \texttt{convert-secret-for-to-static-for}
requires a user-supplied static upper bound.
\item \textbf{Every nonlinearity must be approximated}, and the approximation is
not inferred: \heir{} lists polynomial approximation requiring user annotation
for degree and interval among its stated limitations.
\item \textbf{There is no interactive protocol anywhere in the stack.} No
dialect models secret sharing, garbled circuits, message rounds, or a party.
\end{itemize}

None of these is a defect. They are the correct design consequences of targeting
FHE, where a ciphertext genuinely cannot branch and a comparison genuinely is
not an arithmetic circuit. But together they fix a boundary, and the boundary
has a cost that nobody has measured, because measuring it requires a system on
both sides of it.

\begin{quote}
\textbf{Thesis.} For an important class of programs --- machine-learning
inference above all --- the cheapest and the \emph{most accurate} backend
changes at points \emph{inside} the program, and the switch itself is cheap. A
compiler that can only lower a whole program to one paradigm must pay the worse
option everywhere. Representing the switch requires an execution model in which
``the rest of the computation'' is a first-class object; it is awkward in a
static SSA dataflow IR and natural in a continuation-based one.
\end{quote}

\noindent\textbf{Why ML is the right witness, not an arbitrary application.} A
convolutional network is literally an alternation of the two regimes. Linear
layers are depth-shallow, massively parallel and pack into ciphertext slots ---
the regime FHE was built for. Nonlinearities (ReLU, max-pool, argmax) are
comparisons, which FHE cannot do and must approximate, and which secret sharing
computes exactly in a few rounds. The optimal program for a CNN is therefore
neither an FHE nor an MPC program but an alternation, switching at every layer
boundary --- and the alternation is forced by the network architecture rather
than contrived by us.

% -------------------------------------------------------------------------
\section{The architecture, and what is already built}

\baru{} is a working research compiler and runtime: one Lisp-like source
language over four backends --- plain interpreter, 3-party MPC (ABY3-style
replicated sharing, with a malicious-with-abort mode), FHE (BFV/BGV, TFHE-rs,
and a CKKS arm on OpenFHE), and ZK proofs (Binius64) --- with values crossing
between MPC shares and FHE ciphertexts mid-computation without being exposed.

Its evaluator is a defunctionalized continuation machine: an explicit
\texttt{Continuation} algebraic data type, an \texttt{Eval}/\texttt{Apply} task
trampoline over an explicit stack, and a step budget. Two properties follow, and
we distinguish sharply between what this already delivers and what it enables.

\noindent\textbf{Already delivered.} \emph{(i)} A secret-dependent conditional
is a continuation frame (\texttt{If}, then \texttt{IfResultNormalize}, which
reconciles a numeric result into a share of the right width and owner), so it is
executed under a protocol rather than compiled into a select. \emph{(ii)}
Backend choice is per-operation, by dynamic dispatch through a \texttt{Backend}
trait, rather than fixed by whole-program lowering --- which makes a mid-program
scheme switch an ordinary operation, not an escape hatch. The flagship result
rests on this: the MPC round trip a ReLU already requires \emph{also}
re-randomises the FHE ciphertext, so an arbitrarily deep network runs with
\textbf{zero FHE bootstraps} on an engine supporting one multiplication.

\noindent\textbf{Enabled, and proposed here (Thrust~3).} Continuations are not
yet serialized nor used to suspend at a round boundary --- the transport backend
blocks synchronously. Making the continuation a value turns ``the rest of the
protocol'' into data, permitting \emph{dynamic} round scheduling: $k$
independent subcomputations interleaved on one party, each resumed as its
message arrives. The static scheduler already shows the prize --- exchange count
stayed at 32 for every batch size from 1 to 8 --- but only for statically known
dependency graphs. We propose to generalise it and measure against the scheduler
it subsumes.

% -------------------------------------------------------------------------
\section{What the boundary costs: five measured results}

Full tables in Appendix~A. All figures are measured elapsed times on one machine
with both arms checked against the same plaintext reference.

\medskip
\noindent\textbf{1. The crossover is measured, and it runs both ways.} In
October 2026 we completed the \heir{} baseline described in §4, and it produced
the result this proposal is really about. Both arms evaluate the same MLP on the
same weights with $x^2$ activation --- the same function --- at
\texttt{HEStd\_128\_classic}, on one machine, each verified against the same
plaintext reference.

\begin{center}\small
\begin{tabular}{@{}lrrr@{}}
\toprule
\texttt{64x32x10}, depth 3 & \textbf{\heir{} (CKKS)} & \textbf{2PC, LAN} & \textbf{2PC, WAN} \\
\midrule
one inference & 167.0\,ms & \textbf{14\,ms} & 2113\,ms \\
\bottomrule
\end{tabular}
\end{center}

\noindent\textbf{On LAN secret sharing wins by $11.9\times$; on WAN \heir{}
wins by $12.7\times$} --- a $\mathbf{151\times}$ \textbf{swing from the network
profile alone}, on an identical program with nothing about the cryptography
changed. Appendix~A.1 has the full depth sweep. This is the first time our record
has shown FHE beating the MPC arm, and it is the thesis working rather than
against it: §VI.10--VI.12 measured LeNet-5 at depth 11 with secret sharing ahead
by $32\times$ on latency and $885\times$ past the level budget. Both are true.
\emph{``Is MPC or FHE faster''} is not well-posed, and a compiler that cannot
model the difference cannot make the choice.

\medskip
\noindent\textbf{2. The boundary is a capability cliff, not a crossover in
time.} At multiplicative depth 10 the \heir{} pipeline dies with \textbf{SIGSEGV
during key generation} --- at 554\,MB peak resident against the 564\,MB of the
depth-9 run that succeeds, so \emph{not} an out-of-memory kill, and with no
``cannot find secure parameters'' message either. Depth 9 completes; depth 10
does not run at all, while the secret-sharing curve continues as a straight line
($+7$\,ms per layer on LAN). We have not isolated whether the fault lies in
\heir{}'s parameter selection or in OpenFHE 1.5.1; both are candidates and both
are being reported upstream.

\medskip
\noindent\textbf{3. Approximation returns the wrong answer, at depths that occur.}
This is the result that makes the thesis more than a performance argument. On a
depth sweep of LeNet-5, \textbf{the CKKS arm returned the wrong class at 13
levels} --- the reference top-two margin was 2 quanta and the worst logit error
was 15. The exact (secret-sharing) arm returned the reference class at all five
depths, including margins of one and two quanta. Approximate arithmetic
accumulates error through every rescale and key switch; stochastic truncation's
error is a carry in the last fixed-point bit, unbiased, so it does not
systematically add. Since polynomial approximation of the nonlinearity is the
\emph{only} route available to an FHE-only compiler, this is a correctness
consequence of the architectural boundary, not a tuning issue.

\medskip
\noindent\textbf{4. The performance margin is a choice of regime --- so a
compiler must model the regime.}

\begin{center}\small
\begin{tabular}{@{}lrrl@{}}
\toprule
\textbf{Regime (LeNet-5, same network, same machine)} & \textbf{2-party SS} & \textbf{CKKS} & \textbf{Margin} \\
\midrule
Latency, one inference (50\,ms RTT, 100\,Mbps) & 1.41\,s & 43.8\,s & $32\times$ \\
Throughput, batch 8 (after operand fusion) & 0.363\,s & 2.86\,s & ${\sim}8\times$ \\
Level budget exceeded (forced to bootstrap) & 1.52\,s & 1345.1\,s & $885\times$ \\
\bottomrule
\end{tabular}
\end{center}

\noindent There is also a discontinuity \emph{below} the bootstrap cliff: 13 to
15 levels forces ring $2^{15}\!\to\!2^{16}$ and cost $61.1 \to 160.0$\,s ---
$2.6\times$ for two levels, while multiplication count moves 1.9\%. Security
parameters, not work. Over the same sweep the secret-sharing curve is a straight
line. A line against a staircase: that shape is the compiler's decision surface.

\medskip
\noindent\textbf{5. The language surface determines round complexity.} Writing
the same CNN in \baru{} the way one writes a CNN in a scalar language --- a loop
over sliding windows, comparison per activation --- costs \textbf{507\,920
messages}. Exposing already-existing batched primitives (\texttt{lt\_batch},
\texttt{mul\_batch}, \texttt{a2b\_batch}) to the \emph{language}, then fusing
ReLU with rescale, takes the identical computation to \textbf{33\,106 messages}
and the private diagnosis from 794\,508\,ms to \textbf{911\,ms} --- $872\times$,
with every integer path bit-identical throughout. The frontend, not the
protocol, was the cost. This is the clearest evidence we have that the
language/IR design is the object of study here.

\medskip
\noindent\textbf{What this does not establish.} The CNN margins (result 4) come
from \emph{our own} CKKS arm, $15\times$ off Orion at 3.5\% slot utilisation;
only the MLP rows above are against external code. And the Lisp layer costs real
time --- the same network as a \baru{} program on the plain interpreter takes
299\,ms against 268\,ms for the hand-written Rust pipeline \emph{with image and
model encrypted}. Reducing it is Thrust~3's secondary objective.

% -------------------------------------------------------------------------
\section{The \heir{} comparison, in three tiers}

The comparison must be designed so that the architectural claim is separable
from implementation quality on either side. We propose three tiers, and we
regard Tier~A as a precondition for reporting B or C at all.

\medskip
\noindent\textbf{Tier A --- Parity, as a control condition. \emph{Done for
MLPs; outstanding for CNNs.}} The control is that \baru{} must be within a small
constant factor of \heir{}-generated OpenFHE code at matched security:
\textbf{if our FHE arm is slow, then ``switching to MPC helps'' is an artifact of
a weak baseline rather than a property of the architecture}, and no Tier~B or C
number is attributable. For dense MLPs this is now measured, and §3 reports what
it found --- including \heir{} beating us outright on wide-area MLPs, which is
the control doing its job. Extending it to the convolutional networks of
§VI.10 requires the multi-channel packing of Thrust~1.

\medskip
\noindent\textbf{Tier B --- Expressiveness.} Programs \heir{} cannot express,
each with its mechanism cited rather than asserted: exact ReLU and max-pool (no
comparison in any target scheme); secret-dependent loop bounds
(\texttt{convert-secret-for-to-static-for} requires static bounds); any
mid-program FHE$\leftrightarrow$MPC transition (no interactive dialect). We
report \baru{}'s measured cost and record \heir{} as \emph{not expressible},
naming the pass. A reviewer will object that this compares against a tool not
claiming the capability --- that is the point, stated that way: we quantify what
the exclusion costs on the canonical workload, which needs a system spanning
both sides.

\medskip
\noindent\textbf{Tier C --- Crossover.} Programs whose optimal assignment is
mixed: show the per-region selector finds it and both pure strategies lose. §3
already establishes the crossover exists in both directions; what remains is
showing a compiler can \emph{predict} which side a given program falls on.

\noindent\textbf{Two gates already enforced in code}, because an unfair
comparison is worse than none. \emph{(i)} A public \heir{} issue reports an MNIST
MLP at ${\sim}30$\,s / 3\,GB under OpenFHE's \texttt{HEStd\_NotSet} --- no
security level at all --- against ${\sim}6$\,min / 26\,GB at the secure ring
dimension: a $12\times$ factor between an insecure and a secure baseline. To be
exact, and fair to \heir{}: that insecure run was the reporter's choice, not a
\heir{} default --- \heir{}'s emitted C++ sets only multiplicative depth and
plaintext modulus, leaving OpenFHE's default \texttt{HEStd\_128\_classic} in
force, which we verified by reading the generated code. The gate exists because a
baseline harness must \emph{verify} the parameter it reports rather than assume
it in either direction, and ours refuses to emit a latency unless the run reports
\texttt{HEStd\_\{128,192,256\}\_classic}. \emph{(ii)} Both arms must evaluate one
fixture, with a test asserting the float evaluation reproduces the quantized
classification.

% -------------------------------------------------------------------------
\section{Thrusts, and what the award buys}

\noindent\textbf{T1 --- Multi-channel packing for the FHE arm.} Reach
Orion-class slot utilisation so Tier~A extends to CNNs and all margins are over
the state of the art. Adversarial to our own prior numbers; we will report the
shrunken figures.

\noindent\textbf{T2 --- Extending the \heir{} comparison (where the AWS
credits go).} The baseline now runs, and building it corrected our own estimate
of what it needs: we had deferred it believing secure parameters required
${\sim}26$\,GB, and the secure \texttt{784x32x10} baseline in fact peaks at
532\,MB. So the credits are \emph{not} requested for memory. Two gaps remain
that a single machine cannot close. \emph{(i)} Every wide-area figure we have,
on both arms, comes from a transport that really sleeps rather than a real link
--- and §3 shows the network profile is worth $151\times$, so this is the
largest uncontrolled variable in the whole comparison. Multi-region deployment
(\texttt{sa-east-1} $\leftrightarrow$ \texttt{us-east-1} $\leftrightarrow$
\texttt{eu-west-1}) replaces it with measured intercontinental RTTs.
\emph{(ii)} Separating the depth-10 cliff into ``\heir{} chose bad
parameters,'' ``OpenFHE has a bug,'' and ``this regime needs bootstrapping''
requires the larger instances where secure CKKS bootstrapping is tractable.

\noindent\textbf{T3 --- Serializable continuations and dynamic round scheduling.}
Make the continuation a value; schedule rounds dynamically by resuming whichever
subcomputation's message arrives; measure against the static scheduler it
generalises, on programs whose dependency graph is not statically known.
Secondary objective: close the interpreter overhead reported in §3.

\noindent\textbf{T4 --- The per-region backend selector}, validated by
exhaustive comparison against measured runs of every assignment on small
programs. The research question is not whether a cost model can be written but
\emph{how coarse a model suffices}, with ``suffices'' defined by measurement.

\noindent\textbf{Deliverables.} The open-source tool; a reproducible
matched-threat-model benchmark suite spanning MPC and FHE, which the field
lacks; two publications (compiler/systems on the architecture and backend
selection; applied-crypto on switch-as-bootstrap); \TODO{N} student theses.

% -------------------------------------------------------------------------
\section{Required questions}

\noindent\textbf{Theoretical or applied; domain.} Applied, with a
compiler-architecture contribution and a measurement-methodology contribution.
Domain: compilers and tooling for privacy-preserving computation, evaluated on
private ML inference.

\noindent\textbf{Applications, including for Amazon.} Multi-organization
collaboration on data that cannot be pooled --- an AWS Clean Rooms-shaped
problem, and a listed topic in this call. A cost model spanning MPC and FHE
answers a provisioning question: \emph{given this workload and this network,
which privacy technology, and therefore how much compute?} §3 shows the answer
inverts with the deployment network alone. Secondarily, deciding when
cryptography beats a hardware enclave for a given program region. LGPD-driven
demand for private analytics in Latin America is served from
\texttt{sa-east-1} --- exactly the wide-area shape we measure.

\noindent\textbf{Assumptions and factors that could undermine the results.}
\begin{itemize}
\item \emph{Our own FHE arm is not state of the art} ($15\times$ off Orion,
3.5\% slot utilisation), so every CNN margin in §VI.10 is a margin over our
implementation. The MLP results in §3 do not have this problem --- that arm is
\heir{}-generated OpenFHE, not ours --- which is exactly why Tier~A matters. T1
extends the same footing to CNNs and will reduce our headline margins.
\item \emph{Single-machine provenance.} Wide-area latency is injected by a
sleeping transport, not a real link --- and §3 shows the profile is worth
$151\times$. T2 replaces it; it is the largest uncontrolled variable we have.
\item \emph{The depth/precision sweep uses an untrained fixture}, whose tight
class margins are a fixture property. The \emph{mechanism} is supported; a
quantitative accuracy claim for trained networks needs a trained fixture.
\item \emph{Threat models are matched but not identical in kind.} FHE lets the
client go offline; 2PC needs it present throughout. Matching the adversary model
does not remove this, and it may dominate for some deployments.
\item \emph{Risk to the thesis.} If a well-packed FHE arm wins in every regime
the selector degenerates and the claim weakens to expressiveness alone; the
experiment is designed so that outcome is publishable.
\item \emph{Risk to the comparison.} \heir{} is a moving target; we pin a
release, report its commit, and treat Tier~A as a control, not a scoreboard.
\end{itemize}

\noindent\textbf{Tool: license, materials, maintenance.} License:
\textbf{Apache-2.0} (\texttt{LICENSE}, and declared in \texttt{Cargo.toml}).
Materials: a language reference, a complete experimental record including
negative and withdrawn results (\texttt{docs/experiments.md}), a cross-paradigm
survey with methodology, and a performance handoff written to be read cold.
Maintenance: 129 commits in three months; three contributors, one from USP; 831
tests green; CI on every push.
\emph{Discipline:} every engine \emph{attests} what it implements, and asking
for more is an error, not a silent downgrade. We have exercised it: all ZKP
results predating 2026-07-31 were withdrawn when the shipped verifier proved not
to check constraint satisfaction, and an FHE scheme was deleted rather than
caveated when our own analysis found it insecure at the needed parameters.

% -------------------------------------------------------------------------
\section{Community, students, and budget}

Based at LARC (PCS, Escola Politécnica da USP), which teaches an undergraduate
cryptography course (\emph{Fundamentos de Criptografia}) and supervises graduate
work in applied cryptography; the award supports \TODO{N} MSc/PhD students and
\TODO{N} undergraduate scholars. A compiler project absorbs students unusually
well --- backends, cost models, gadget libraries and benchmark suites
parallelize into defensible individual theses with measurable results --- and
the matched-threat-model benchmark suite is itself a community contribution
whose absence caused a comparison we had to withdraw. Applied-cryptography
funding concentrates in North America, Europe and Israel; USP is Latin America's
leading research university, and this supports a group training Brazilian
students in a subfield where the region is under-represented.

\medskip
\noindent\textbf{Budget} (unrestricted, excluding overhead; \TODO{confirm
amounts against USP/FUSP scholarship rates}). Graduate student support
(\TODO{N}, 12 mo.) \TODO{amt}; undergraduate \emph{iniciação científica}
scholarships (\TODO{N}, 12 mo.) \TODO{amt}; conference travel, open-access
charges and local compute \TODO{amt}. \textbf{Total unrestricted
\TODO{$\leq$ 80{,}000} USD.} \textbf{AWS promotional credits}
(\TODO{$\leq$ 40{,}000}) are requested for Thrust~2 specifically, per §5.

% -------------------------------------------------------------------------
\vspace{0.6ex}\hrule\vspace{0.6ex}
{\footnotesize\noindent\textbf{References.}
[1]~\heir{} Contributors. \emph{HEIR: A Compiler Infrastructure for Homomorphic
Encryption}. arXiv:2508.11095, 2025.\quad
[2]~A.~Ebel, G.~Garimella, B.~Reagen. \emph{Orion: A Fully Homomorphic
Encryption Framework for Deep Learning}. ASPLOS, 2025.\par}

\clearpage
% =========================================================================
\appendix
\section*{Appendix A --- Measurement detail}
\setcounter{section}{1}

\textbf{Provenance.} Every figure is an elapsed wall-clock time or an
instrumented counter from a real run on one machine, recorded in
\texttt{docs/experiments.md} alongside the command that produced it. No figure
is composed from a cost model or extrapolated. Where a result has not been
measured it is described as outstanding rather than estimated.

\subsection{\heir{} vs.\ 2-party MPC, same function, same machine}

OpenFHE 1.5.1 from source; \texttt{heir-py} 2026.8.1; every row reports
\texttt{HEStd\_128\_classic} and reproduced the plaintext classification, so
both §4 gates passed. Activation $x^2$ on both arms. \heir{} \texttt{eval}
excludes keygen; MPC \texttt{online} excludes preprocessing. MPC rows are the
\texttt{dealer} \emph{lower bound}; deployable \texttt{lowgear} totals are given
below. LAN is 0\,ms / 10\,Gbit/s, WAN 50\,ms one-way / 100\,Mbit/s.

\begin{center}\small
\begin{tabular}{@{}lrrrrrr@{}}
\toprule
\textbf{Network} & \textbf{Depth} & \textbf{\heir{} eval} & \textbf{keygen} & \textbf{MPC LAN} & \textbf{MPC WAN} & \textbf{Rounds} \\
\midrule
\texttt{64x32x10}            & 3  & 167.0\,ms & 91.7\,ms  & 14\,ms & 2113\,ms & 41 \\
\texttt{784x32x10}           & 3  & 178.0\,ms & 105.3\,ms & 26\,ms & 2167\,ms & 41 \\
\texttt{64x32x32x10}         & 5  & 659.3\,ms & 271.2\,ms & 20\,ms & 3044\,ms & 59 \\
\texttt{64x32x32x32x10}      & 7  & 1051.3\,ms & 347.6\,ms & 27\,ms & 3978\,ms & 77 \\
\texttt{64x32x32x32x32x10}   & 9  & 1493.4\,ms & 436.3\,ms & 34\,ms & 4909\,ms & 95 \\
\texttt{64x32($\times$5)x10} & 10 & \textbf{SIGSEGV} & --- & 41\,ms & 5842\,ms & 113 \\
\bottomrule
\end{tabular}
\end{center}

\noindent Deployable \texttt{lowgear} totals including preprocessing:
\texttt{64x32x10} 264\,ms LAN / 3264\,ms WAN; \texttt{784x32x10} 388\,ms LAN /
3791\,ms WAN. Measured slopes per added hidden layer: \heir{}
${\approx}440$\,ms, MPC ${\approx}930$\,ms on WAN but ${\approx}7$\,ms on LAN.

\medskip
\noindent\textbf{Four things this does not say.} \emph{(i)} These are dense
MLPs, not the convolutional LeNet-5 of §A.2--A.3; the two sets compare only
within themselves. \emph{(ii)} \texttt{dealer} is a lower bound --- against
\texttt{lowgear} totals the LAN comparison is roughly a tie
(264\,ms vs.\ $167+92$\,ms), and the defensible LAN claim is the \emph{online}
one. \emph{(iii)} The interaction asymmetry is unchanged: \heir{}'s client can
go offline, ours cannot. \emph{(iv)} Our arm computes the argmax inside the
protocol and reveals only the class, while \heir{} returns the full logit
vector --- strictly more work, strictly less disclosure.

\subsection{Depth sweep: a line against a staircase}

LeNet-5 with additional dense-plus-activation blocks (two levels each), one
variable.

\begin{center}\small
\begin{tabular}{@{}rrrrrrr@{}}
\toprule
\textbf{Levels} & \textbf{Exch} & \textbf{MiB} & \textbf{MPC, 50\,ms WAN} & \textbf{Ring} & \textbf{CKKS} & \textbf{Ratio} \\
\midrule
11 & 32 & 4.82 & 1517\,ms & $2^{15}$ & 43.8\,s & $29\times$ \\
13 & 36 & 4.96 & 1683\,ms & $2^{15}$ & 61.1\,s & $36\times$ \\
15 & 40 & 5.10 & 1804\,ms & $2^{16}$ & 160.0\,s & $89\times$ \\
17 & 44 & 5.23 & 1937\,ms & --- & --- & \\
19 & 48 & 5.37 & 2110\,ms & --- & --- & \\
\bottomrule
\end{tabular}
\end{center}

\subsection{Precision: where approximation returns the wrong class}

Margins in units of the fixed-point quantum ($2^{-8}$).

\begin{center}\small
\begin{tabular}{@{}rrrrll@{}}
\toprule
\textbf{Levels} & \textbf{Ref.\ class} & \textbf{Margin} & \textbf{CKKS worst error} & \textbf{CKKS verdict} & \textbf{MPC} \\
\midrule
11 & 6 & 58\,q & 16\,q & agrees & 6 \\
13 & 0 & \textbf{2\,q} & 15\,q & \textbf{wrong class (7)} & \textbf{0} \\
15 & 8 & 10\,q & 21\,q & agrees, fortuitously & 8 \\
17 & 2 & \textbf{1\,q} & --- & --- & \textbf{2} \\
19 & 3 & 72\,q & --- & --- & 3 \\
\bottomrule
\end{tabular}
\end{center}

\subsection{The frontend determines round complexity ($872\times$)}

MedMNIST pneumonia, one diagnosis, same weights, every integer path bit-identical.

\begin{center}\small
\begin{tabular}{@{}lrrr@{}}
\toprule
\textbf{Version} & \textbf{Plaintext, 10 img} & \textbf{Private (image)} & \textbf{Messages} \\
\midrule
Scalar Lisp & 3465\,ms & 794\,508\,ms & 507\,920 \\
+ batched vector primitives & 322\,ms & 6446\,ms & 33\,181 \\
+ fused \texttt{vrelu} & 299\,ms & 4326\,ms & 33\,154 \\
+ garbled nonlinearity & --- & \textbf{911\,ms} & 33\,106 \\
\midrule
\textbf{Improvement} & $\mathbf{11.6\times}$ & $\mathbf{872\times}$ & $\mathbf{15.4\times}$ \\
\bottomrule
\end{tabular}
\end{center}

\subsection{Backend-aware transformations, each selected by a probe}

\begin{center}\small
\begin{tabular}{@{}llr@{}}
\toprule
\textbf{Transformation} & \textbf{Why it is a backend-aware decision} & \textbf{Measured} \\
\midrule
Argmax as tournament, not fold & same comparisons, $n{-}1\to\lceil\log_2 n\rceil$ depth & $1.67\times$ \\
Rescale by truncation pair & query-independent; crosses the online/offline line & $3.67\times$ \\
All-pairs argmax & spends bandwidth headroom to buy constant depth & $1.47\times$ \\
Open each operand once & reuse appears only at batch $>1$ & $10.8\times$ traffic \\
\midrule
\multicolumn{2}{@{}l}{\textbf{LeNet-5, wide-area, cumulative}} & $\mathbf{7.76\,s\to1.41\,s}$ \\
\bottomrule
\end{tabular}
\end{center}

\noindent All-pairs argmax is the thesis in miniature: it performs 45
comparisons instead of 9, and is correct \emph{only} because a separate
measurement showed that workload had bandwidth headroom. On a bandwidth-bound
link the opposite lowering wins, so both are kept and selected by cost model ---
a decision a single-paradigm compiler has no basis to make.

\section*{Appendix B --- Capability matrix}
\setcounter{section}{2}

Rows are sourced to \heir{}'s paper (arXiv:2508.11095); none is inferred.

\begin{center}\small
\begin{tabular}{@{}lll@{}}
\toprule
\textbf{Capability} & \textbf{\heir{}} & \textbf{\baru{}} \\
\midrule
FHE schemes & BGV/BFV, CKKS, CGGI & BFV/BGV, CKKS (OpenFHE), TFHE-rs \\
Backend libraries & OpenFHE, Lattigo, TFHE-rs, Jaxite & OpenFHE; own NTT \\
Hardware targets & several (\texttt{pisa}, accelerators) & none \\
Secret sharing / MPC & not present & ABY3; malicious-with-abort \\
Garbled circuits & not present & yes (nonlinearity lowering) \\
FHE $\leftrightarrow$ MPC switch & not present & both directions \\
Exact ReLU / max-pool & not expressible (no comparison) & exact \\
Nonlinearity route & polynomial approx., user-annotated & exact, or approx.\ by choice \\
Secret-dependent branch & \texttt{convert-if-to-select} (both arms) & executed under protocol \\
Secret-dependent loop bound & unsupported (static bound required) & executed under protocol \\
BGV/BFV bootstrapping & unsupported (no backend) & avoided via switch-as-bootstrap \\
Measured depth ceiling & SIGSEGV at mult.\ depth 10 & no depth budget; linear in depth \\
Maturity, ecosystem, adoption & \textbf{substantially ahead} & research prototype \\
\bottomrule
\end{tabular}
\end{center}

\noindent The last row is not a formality. \heir{} is production-grade
infrastructure with institutional backing, multiple hardware targets and a
community; \baru{} is a research prototype by a small group. The claim in this
proposal is confined to the rows above it, and specifically to what the
architecture makes \emph{expressible} --- which is why Tier~A parity is a
precondition rather than a result.

\end{document}
